Thus far the calculi studied have only covered the implicational
fragments of the two approaches, i.e., just function types and
variables.  Adding product and sum types to the calculi brings the
theory closer to a useful basis for understanding actual
programming languages such as the linear-types extension to Haskell (graded-base) or Granule
(linear-base).
We thus now add products and sums to Linear Base and Graded Modal Base
following standard approaches in the literature for both flavours
of linear-base and graded-base systems.  This raises interesting questions
that delineate the expressive power of linear- vs. graded-base
calculi. From linear-base to graded-base, a straightforward
homomorphic translation is possible. However, for the opposite
direction we will see a difference in expressivity that needs handling
by extending the linear-base system.

\paragraph{Syntax} We extend Linear Base and Graded Modal Base (which already share the same
syntax) to \textbf{Linear~Core} and \textbf{Graded~Modal~Core} respectively, with the same extended syntax:
%
\begin{align*}
  \SYSTEMnt{t} & ::= ... \mid  \langle  \SYSTEMnt{t_{{\mathrm{1}}}} ,  \SYSTEMnt{t_{{\mathrm{2}}}}  \rangle  \mid  \mathsf{let} \, \langle  \SYSTEMmv{x} ,  \SYSTEMmv{y}  \rangle =  \SYSTEMnt{t_{{\mathrm{1}}}}  \, \mathsf{in} \,  \SYSTEMnt{t_{{\mathrm{2}}}}  \mid  \langle \rangle  \mid  \mathsf{let} \, \langle \rangle =  \SYSTEMnt{t_{{\mathrm{1}}}}  \, \mathsf{in} \,  \SYSTEMnt{t_{{\mathrm{2}}}}  \\
         \mid & \; \mathsf{inj}_1 \,  \SYSTEMnt{t}  \mid  \mathsf{inj}_2 \,  \SYSTEMnt{t}  \mid  \mathsf{case} \,  \SYSTEMnt{t}  \, \mathsf{of} \, \{ \mathsf{inj1} \,  \SYSTEMmv{x}  \rightarrow  \SYSTEMnt{t_{{\mathrm{1}}}}  ; \, \mathsf{inj2} \,  \SYSTEMmv{y}  \rightarrow  \SYSTEMnt{t_{{\mathrm{2}}}}  \} 
  \tag{terms}
\end{align*}
%
These additional term formers provide product introduction and elimination, unit introduction
and elimination, and sum (coproduct) introductions and elimination respectively.

The operational semantics are the standard call-by-name operational semantics for products and sums given in Figure~\ref{fig:linear-core} for Linear Core, with the same
added to operational semantics of Graded Modal Core. We elide the $\beta\eta$-equational theory which is standard (see Appendices~\refappendix{app:definitions-linear-core} and
~\refappendix{app:definitions-graded-core}).

\begin{figure}[t]
\raggedright\noindent\framebox{$ \SYSTEMnt{t}  \rightsquigarrow_{\textsc{l} }  \SYSTEMnt{t} $} \quad Standard operational semantics (Call-By-Name) for products, units, and sums
\begin{gather*}
  \begin{array}{c}
  \SYSTEMdruleSemLinprodCong{}
  \;
  \SYSTEMdruleSemLinprodBeta{}
  \\[1.5em]
  \SYSTEMdruleSemLinunitCong{}
  \quad
  \SYSTEMdruleSemLinunitBeta{}
  \\[1.5em]
  \SYSTEMdruleSemLincongCase{}
  \\[1.5em]
  \SYSTEMdruleSemLincaseInjOne{}
  \\[1.5em]
  \SYSTEMdruleSemLincaseInjTwo{}
  \end{array}
\end{gather*}
\bigskip

\raggedright\noindent\framebox{$ \Gamma  \vdash_{\textsc{l} }  \SYSTEMnt{t}  :  \SYSTEMnt{A} $} \quad Typing rules
\begin{gather*}
\begin{align*}
\begin{array}{c}
\SYSTEMdruleLinprodi{}
\quad
\SYSTEMdruleLinprode{}
\\[1.5em]
\SYSTEMdruleLinuniti{}
\quad
\SYSTEMdruleLinunite{}
\quad
\SYSTEMdruleLinsumiOne{}
\\[1.5em]
\SYSTEMdruleLinsumiTwo{}
\quad
\SYSTEMdruleLinsume{}
\end{array}
\end{align*}
\end{gather*}
\caption{Linear Core semantics and typing}
\label{fig:linear-core}
\end{figure}

% The additional syntax for types differs only in
%the symbol used for the type constructors of sums and products.
% ($     \otimes     $ vs $     \times     $ and $     \oplus     $ vs $     +     $).
% \newpage
To type the new syntax, we extend the type system of Linear Base to Linear Core
(Section~\ref{sec:linear-core}), and extend Graded Modal Base to
Graded Modal Core (Section~\ref{sec:graded-core}).
Section~\ref{sec:linear-core-to-graded} then gives the
translation from Linear Core to Graded Modal Core which is
straightforwardly homomorphic.
Section~\ref{sec:no-trans-graded-core-to-linear-core} then
shows that the opposite direction is not possible, requiring an
extension of Linear Core to \textbf{Linear~Push~Core} in
Section~\ref{sec:linear-push-core}, at the cost of weakening the correspondence between our $\Box$-modality
and Linear Logic's $!$-modality.  We then give the translation
from Graded Modal Core to Linear Push Core
(Section~\ref{sec:trans-graded-core-to-linear-push-core}) and the translation
from Linear Push Core to Graded Modal Core
(Section~\ref{sec:trans-linear-push-core-to-graded-core}).

\subsection{Linear Core}
\label{sec:linear-core}

Linear Base augmented with products and sums becomes what we call \textbf{Linear~Core}.
Linear Core features the following new type constructors for products,
units, and sums:
%
\begin{align*}
  \SYSTEMnt{A} & ::= ... \mid  \SYSTEMnt{A}  \otimes  \SYSTEMnt{B}  \mid  \mathrm{unit}  \mid  \SYSTEMnt{A}  \oplus  \SYSTEMnt{B} 
  \tag{types}
\end{align*}

\paragraph{Typing} Linear Core extends Linear Base typing (given in Figure~\ref{fig:linear-base-types})
with the judgments for multiplicative products and their unit, along with additive disjunction, in Figure~\ref{fig:linear-core}, which
are standard from Linear Logic~\cite{girard1987linear}. This is the same typing as in linear-base systems such as Granule~\cite{DBLP:journals/pacmpl/OrchardLE19,DBLP:journals/corr/abs-2112-14966}, the synthesis calculus of \citet{hughes2021linear}, and that of \citet{brunel2014core} (although
their system has a less general natural numbers data type instead of products and sums, but the same insights apply).

As in Linear Base, any terms comprising multiple sub-terms have their contexts
combined using context addition $+$ which adds together the grades of
graded assumptions using semiring addition.



\subsection{Graded Modal Core}
\label{sec:graded-core}

Graded Modal Base augmented with products and sums becomes \textbf{Graded Modal Core}.
Graded Modal Core features the following new type constructors for products, units, and sums:
%
\begin{align*}
  \SYSTEMnt{A} & ::= ... \mid  \SYSTEMnt{A}  \times  \SYSTEMnt{B}  \mid  \mathrm{unit}  \mid  \SYSTEMnt{A}  +  \SYSTEMnt{B} 
  \tag{types}
\end{align*}

\paragraph{Typing}

Graded Modal Core extends Graded Modal Base typing (given in Section~\ref{sec:graded-modal-base}) by the judgments
given in Figure~\ref{fig:graded-base-types-products-sums}. The typing is akin to that shown
by \citet{DBLP:journals/pacmpl/AbelB20}, in Linear Haskell~\cite{DBLP:journals/pacmpl/BernardyBNJS18}, and by \citet{grad} (although they restrict their
products initially to $r = 1$).

In the elimination of the products we may use each variable
according to grade $\SYSTEMnt{r}$ and thus the context $\Delta_{{\mathrm{1}}}$ for the eliminated product
is multiplied by $\SYSTEMnt{r}$ in the conclusion. Thus, the usage of the components is propagated
to demands on the dependencies for the product. The elimination of sums
has a similar idea, propagating demands on the variables $x$ and $y$ through to the
dependencies of the eliminated sum term $t$. There is an additional side condition here too
 ($ \SYSTEMsym{1}  \, \textcolor{coeffectColor}{\sqsubseteq} \,  \SYSTEMnt{r} $) accounting for usage incurred by the inspection of the sum
constructor, and thus some information gained. This side condition is key to applications of
graded coeffect types for information-flow control~(\ref{exm:security-lattice}) \cite{DBLP:journals/pacmpl/AbelB20,mycroftfest2024}:
if the grade here is $r = \mathsf{Hi}$, representing
a high-security value, then the side condition fails as $1 = \mathsf{Lo}$ and
thus $\mathsf{Lo} \not\succeq \mathsf{Hi}$. This prevents pattern matching on a high-security
value as a control-flow attack to leak information. Lastly note the elimination of
the unit type similarly allows (an arbitrary) scaling of the scrutinee's context, in line with
presentations in the literature (e.g.,~\cite{DBLP:journals/pacmpl/AbelB20}).
%
\begin{figure}[t]
\raggedright\noindent\framebox{$ \Delta  \vdash_{\textsc{g} }  \SYSTEMnt{t}  :  \SYSTEMnt{A} $} \quad Typing rules
\begin{gather*}
\begin{align*}
\begin{array}{c}
\SYSTEMdruleGradprodi{}
\quad
\SYSTEMdruleGradprode{}
\\[1.25em]
\SYSTEMdruleGraduniti{}
\quad
\SYSTEMdruleGradunite{}
\quad
\SYSTEMdruleGradsumiOne{}
\\[1.25em]
\SYSTEMdruleGradsumiTwo{}
\;
\SYSTEMdruleGradsume{}
\end{array}
\end{align*}
\end{gather*}
\caption{Graded Modal Core additional typing rules.}
\label{fig:graded-base-types-products-sums}
\end{figure}

\subsection{Translation from Linear Core to Graded Modal Core}
\label{sec:linear-core-to-graded}

We extend the translation from Section \ref{subsec:transl-lb2gmb} via a homomorphism
on the additional syntax:
%
% TEMPORARILY REDEFINE SYSTEM DRULE FOR THE TRANSLATION WHICH IS REALLY
% JUST A FUNCTIOn
\renewcommand{\triangleq}{& \oldTriangleq &}
\renewcommand{\equiv}{& \oldEquiv &}
\renewcommand{\SYSTEMdrule}[4][]{#3 \qquad #2}
%
\begin{gather*}
{\setlength{\arraycolsep}{0.1em}
\begin{array}{lcl@{\hspace{0em}}r}
\SYSTEMdruleTranslLtoGTmprodi{} \\
\SYSTEMdruleTranslLtoGTmprode{} \\
\SYSTEMdruleTranslLtoGTmuniti{} \\
\SYSTEMdruleTranslLtoGTmunite{} \\
\SYSTEMdruleTranslLtoGTmsumiOne{} \\
\SYSTEMdruleTranslLtoGTmsumiTwo{} & \text{(Term translation)} \\
\SYSTEMdruleTranslLtoGTmsume{}  \\[0.5em]
\SYSTEMdruleTranslLtoGTyprodTy{} \\
\SYSTEMdruleTranslLtoGTysumTy{} \\
\SYSTEMdruleTranslLtoGTyunitTy{} & \text{(Type translation)}
\end{array}}
\end{gather*}
%
\let\SYSTEMdrule\oldSYSTEMdrule
\let\equiv\oldEquiv
\let\triangleq\oldTriangleq
%
We can then re-establish the two key results about the translation (proofs in
Appendix~\refappendix{app:proofs-lin-core-to-grad-core}):
%
\begin{restatable}{thrm}{linCanToGradTranslationExt}[Soundness of the Linear Core to Graded Modal Core translation]
 \begin{itemize}[itemsep=0.2em]
  \item Type preservation: $ \Gamma  \vdash_{\textsc{l} }  \SYSTEMnt{t}  :  \SYSTEMnt{A}  \implies   \textcolor{LGcolor}{\llparenthesis}  \Gamma  \textcolor{LGcolor}{\rrparenthesis}   \vdash_{\textsc{g} }   \textcolor{LGcolor}{\llparenthesis} \smidge  \SYSTEMnt{t}  \smidge \textcolor{LGcolor}{\rrparenthesis}   :   \textcolor{LGcolor}{\llparenthesis} \smidge  \SYSTEMnt{A}  \smidge \textcolor{LGcolor}{\rrparenthesis}  $;
  \item Operational correspondence: $ \SYSTEMnt{t}  \rightsquigarrow_{\textsc{l} }  \SYSTEMnt{t'}  \implies   \textcolor{LGcolor}{\llparenthesis} \smidge  \SYSTEMnt{t}  \smidge \textcolor{LGcolor}{\rrparenthesis}   \rightsquigarrow_{\textsc{g} }   \textcolor{LGcolor}{\llparenthesis} \smidge  \SYSTEMnt{t'}  \smidge \textcolor{LGcolor}{\rrparenthesis}  $;
  \item Equation preservation: $ \SYSTEMnt{t}  \,\equiv_{\textsc{l} }\,  \SYSTEMnt{t'}  \implies   \textcolor{LGcolor}{\llparenthesis} \smidge  \SYSTEMnt{t}  \smidge \textcolor{LGcolor}{\rrparenthesis}   \equiv_{\textsc{g} }   \textcolor{LGcolor}{\llparenthesis} \smidge  \SYSTEMnt{t'}  \smidge \textcolor{LGcolor}{\rrparenthesis}  $.
\end{itemize}
\end{restatable}
%
A notable point of the type preservation proof is that when translating sum elimination from Linear Core to
Graded Modal Core, we pick the grade $r = 1$ which satisfies the side condition that $ \SYSTEMsym{1}  \, \textcolor{coeffectColor}{\sqsubseteq} \,  \SYSTEMsym{1} $
trivially by reflexivity of the pre-order.
%

\subsection{No Type-Sound Translation from Graded Modal Core to Linear Core}
\label{sec:no-trans-graded-core-to-linear-core}

We now attempt a similar homomorphic translation ($  \textcolor{GLcolor}{\llbracket}   -   \textcolor{GLcolor}{\rrbracket}  _\circ $) back from Graded Modal Core
to Linear Core, focusing on just products and sums:
%
\renewcommand{\equiv}{ \oldEquiv}
\renewcommand{\SYSTEMdrule}[4][]{{\displaystyle{#3}}}
%
\begin{align}
 \tag{Term translation} \\[-1.5em]
\notag & \begin{array}{l}
  \textcolor{GLcolor}{\llbracket}   \langle  \SYSTEMnt{t_{{\mathrm{1}}}} ,  \SYSTEMnt{t_{{\mathrm{2}}}}  \rangle   \textcolor{GLcolor}{\rrbracket}  _\circ  \triangleq  \langle    \textcolor{GLcolor}{\llbracket}  \SYSTEMnt{t_{{\mathrm{1}}}}  \textcolor{GLcolor}{\rrbracket}  _\circ  ,    \textcolor{GLcolor}{\llbracket}  \SYSTEMnt{t_{{\mathrm{2}}}}  \textcolor{GLcolor}{\rrbracket}  _\circ   \rangle  \\
  \textcolor{GLcolor}{\llbracket}   \mathsf{let} \, \langle  \SYSTEMmv{x} ,  \SYSTEMmv{y}  \rangle =  \SYSTEMnt{t_{{\mathrm{1}}}}  \, \mathsf{in} \,  \SYSTEMnt{t_{{\mathrm{2}}}}   \textcolor{GLcolor}{\rrbracket}  _\circ  \triangleq   \mathsf{let} \, \langle  \SYSTEMmv{x} ,  \SYSTEMmv{y}  \rangle =    \textcolor{GLcolor}{\llbracket}  \SYSTEMnt{t_{{\mathrm{1}}}}  \textcolor{GLcolor}{\rrbracket}  _\circ   \, \mathsf{in} \,   \textcolor{GLcolor}{\llbracket}  \SYSTEMnt{t_{{\mathrm{2}}}}  \textcolor{GLcolor}{\rrbracket}   _\circ \\
  \textcolor{GLcolor}{\llbracket}   \mathsf{inj}_i \,  \SYSTEMnt{t}   \textcolor{GLcolor}{\rrbracket}  _\circ  \triangleq   \mathsf{inj}_i \,   \textcolor{GLcolor}{\llbracket}  \SYSTEMnt{t}  \textcolor{GLcolor}{\rrbracket}   _\circ \\
  \textcolor{GLcolor}{\llbracket}   \mathsf{inj}_2 \,  \SYSTEMnt{t}   \textcolor{GLcolor}{\rrbracket}  _\circ  \triangleq   \mathsf{inj}_2 \,   \textcolor{GLcolor}{\llbracket}  \SYSTEMnt{t}  \textcolor{GLcolor}{\rrbracket}   _\circ \\
  \textcolor{GLcolor}{\llbracket}   \mathsf{case} \,  \SYSTEMnt{t}  \, \mathsf{of} \, \{ \mathsf{inj1} \,  \SYSTEMmv{x}  \rightarrow  \SYSTEMnt{t_{{\mathrm{1}}}}  ; \, \mathsf{inj2} \,  \SYSTEMmv{y}  \rightarrow  \SYSTEMnt{t_{{\mathrm{2}}}}  \}   \textcolor{GLcolor}{\rrbracket}  _\circ  \triangleq  \mathsf{case} \,    \textcolor{GLcolor}{\llbracket}  \SYSTEMnt{t}  \textcolor{GLcolor}{\rrbracket}  _\circ   \, \mathsf{of} \, \{ \mathsf{inj1} \,  \SYSTEMmv{x}  \rightarrow    \textcolor{GLcolor}{\llbracket}  \SYSTEMnt{t_{{\mathrm{1}}}}  \textcolor{GLcolor}{\rrbracket}  _\circ   ; \, \mathsf{inj2} \,  \SYSTEMmv{y}  \rightarrow    \textcolor{GLcolor}{\llbracket}  \SYSTEMnt{t_{{\mathrm{2}}}}  \textcolor{GLcolor}{\rrbracket}  _\circ   \} 
\end{array} \\
\tag{Type translation}
& \begin{array}{l}
  \textcolor{GLcolor}{\llbracket}   \SYSTEMnt{A}  \times  \SYSTEMnt{B}   \textcolor{GLcolor}{\rrbracket}  _\circ  \triangleq     \textcolor{GLcolor}{\llbracket}  \SYSTEMnt{A}  \textcolor{GLcolor}{\rrbracket}  _\circ   \otimes   \textcolor{GLcolor}{\llbracket}  \SYSTEMnt{B}  \textcolor{GLcolor}{\rrbracket}   _\circ  \\
  \textcolor{GLcolor}{\llbracket}   \SYSTEMnt{A}  +  \SYSTEMnt{B}   \textcolor{GLcolor}{\rrbracket}  _\circ  \triangleq     \textcolor{GLcolor}{\llbracket}  \SYSTEMnt{A}  \textcolor{GLcolor}{\rrbracket}  _\circ   \oplus   \textcolor{GLcolor}{\llbracket}  \SYSTEMnt{B}  \textcolor{GLcolor}{\rrbracket}   _\circ 
\end{array}
\end{align}
%
\let\SYSTEMdrule\oldSYSTEMdrule
\let\equiv\oldEquiv
%
However, this translation is unsound: consider the following Graded Modal Core terms, $f$ and $g$,
which manifest a distributivity of grading over products and sums:
%
\begin{align*}
\begin{array}{rl}
f &:   (   \SYSTEMnt{A}  \times  \SYSTEMnt{B}   )   \xrightarrow{\textcolor{coeffectColor}{ \SYSTEMnt{r} } }   (    \textcolor{coeffectColor}{\square_{ \SYSTEMnt{r} } }  \SYSTEMnt{A}   \times   \textcolor{coeffectColor}{\square_{ \SYSTEMnt{r} } }  \SYSTEMnt{B}    )   \\
f &=  \lambda  \SYSTEMmv{p}  .  \SYSTEMsym{(}   \mathsf{let} \, \langle  \SYSTEMmv{x} ,  \SYSTEMmv{y}  \rangle =  \SYSTEMmv{p}  \, \mathsf{in} \,   \langle   \textcolor{coeffectColor}{[}  \SYSTEMmv{x}  \textcolor{coeffectColor}{]}  ,   \textcolor{coeffectColor}{[}  \SYSTEMmv{y}  \textcolor{coeffectColor}{]}   \rangle    \SYSTEMsym{)} 
\end{array}
\begin{array}{rl}
g &:   (   \SYSTEMnt{A}  +  \SYSTEMnt{B}   )   \xrightarrow{\textcolor{coeffectColor}{ \SYSTEMnt{r} } }   (    \textcolor{coeffectColor}{\square_{ \SYSTEMnt{r} } }  \SYSTEMnt{A}   +   \textcolor{coeffectColor}{\square_{ \SYSTEMnt{r} } }  \SYSTEMnt{B}    )   \quad (\text{where }  \SYSTEMsym{1}  \, \textcolor{coeffectColor}{\sqsubseteq} \,  \SYSTEMnt{r} )\\
g &=  \lambda  \SYSTEMmv{e}  .   \mathsf{case} \,  \SYSTEMmv{e}  \, \mathsf{of} \, \{ \mathsf{inj1} \,  \SYSTEMmv{x}  \rightarrow   \mathsf{inj}_1 \,   \textcolor{coeffectColor}{[}  \SYSTEMmv{x}  \textcolor{coeffectColor}{]}    ; \, \mathsf{inj2} \,  \SYSTEMmv{y}  \rightarrow   \mathsf{inj}_2 \,   \textcolor{coeffectColor}{[}  \SYSTEMmv{y}  \textcolor{coeffectColor}{]}    \}  
\end{array}
\end{align*}
%
Thus $f$ is a function that distributes the resource requirements on a pair onto its components, and similarly, $g$ distributes the resource requirements on a sum onto the components of its injections.

Applying the above translation to $f$ and $g$ yields the following ill-typed terms
(to avoid clutter, we consider $  \textcolor{GLcolor}{\llbracket}  \SYSTEMnt{A}  \textcolor{GLcolor}{\rrbracket}  _\circ  $ interchangeable with $ \SYSTEMnt{A} $):
%
\begin{align*}
  \textcolor{GLcolor}{\llbracket}  \SYSTEMmv{f}  \textcolor{GLcolor}{\rrbracket}  _\circ  &:   \textcolor{coeffectColor}{\square_{ \SYSTEMnt{r} } }   (   \SYSTEMnt{A}  \otimes  \SYSTEMnt{B}   )    \multimap   (    \textcolor{coeffectColor}{\square_{ \SYSTEMnt{r} } }  \SYSTEMnt{A}   \otimes   \textcolor{coeffectColor}{\square_{ \SYSTEMnt{r} } }  \SYSTEMnt{B}    )   \tag{ill-typed} \\
\notag
  \textcolor{GLcolor}{\llbracket}  \SYSTEMmv{f}  \textcolor{GLcolor}{\rrbracket}  _\circ  &=  \lambda  \SYSTEMmv{b}  .  \SYSTEMsym{(}   \mathsf{let} \, \textcolor{coeffectColor}{[}  \SYSTEMmv{p}  \textcolor{coeffectColor}{]} =  \SYSTEMmv{b}  \, \mathsf{in} \,   \mathsf{let} \, \langle  \SYSTEMmv{x} ,  \SYSTEMmv{y}  \rangle =  \SYSTEMmv{p}  \, \mathsf{in} \,   \langle   \textcolor{coeffectColor}{[}  \SYSTEMmv{x}  \textcolor{coeffectColor}{]}  ,   \textcolor{coeffectColor}{[}  \SYSTEMmv{y}  \textcolor{coeffectColor}{]}   \rangle     \SYSTEMsym{)}  \\[.25em]
\notag
  \textcolor{GLcolor}{\llbracket}  \SYSTEMmv{g}  \textcolor{GLcolor}{\rrbracket}  _\circ  &:   \textcolor{coeffectColor}{\square_{ \SYSTEMnt{r} } }   (   \SYSTEMnt{A}  \oplus  \SYSTEMnt{B}   )    \multimap   (    \textcolor{coeffectColor}{\square_{ \SYSTEMnt{r} } }  \SYSTEMnt{A}   \oplus   \textcolor{coeffectColor}{\square_{ \SYSTEMnt{r} } }  \SYSTEMnt{B}    )   \\
\notag
  \textcolor{GLcolor}{\llbracket}  \SYSTEMmv{g}  \textcolor{GLcolor}{\rrbracket}  _\circ  &=  \lambda  \SYSTEMmv{b}  .  \SYSTEMsym{(}   \mathsf{let} \, \textcolor{coeffectColor}{[}  \SYSTEMmv{e}  \textcolor{coeffectColor}{]} =  \SYSTEMmv{b}  \, \mathsf{in} \,   \mathsf{case} \,  \SYSTEMmv{e}  \, \mathsf{of} \, \{ \mathsf{inj1} \,  \SYSTEMmv{x}  \rightarrow   \mathsf{inj}_1 \,   \textcolor{coeffectColor}{[}  \SYSTEMmv{x}  \textcolor{coeffectColor}{]}    ; \, \mathsf{inj2} \,  \SYSTEMmv{y}  \rightarrow   \mathsf{inj}_2 \,   \textcolor{coeffectColor}{[}  \SYSTEMmv{y}  \textcolor{coeffectColor}{]}    \}    \SYSTEMsym{)} 
\end{align*}
%
To see why these are ill-typed, the following gives a partial derivation for the body of
$  \textcolor{GLcolor}{\llbracket}  \SYSTEMmv{f}  \textcolor{GLcolor}{\rrbracket}  _\circ $:

\begin{gather*}
    \inferrule*[right=\SYSTEMRenameRuleLinlet{}]
      {
       % \inferrule*[Right=\SYSTEMRenameRuleLinvar{}, vdots=5.5em, rightskip=12em]
      %    { }
       %   {   \SYSTEMmv{b}  :   \textcolor{coeffectColor}{\square_{ \SYSTEMnt{r} } }   (   \SYSTEMnt{A}  \otimes  \SYSTEMnt{B}   )      \vdash_{\textsc{l} }  \SYSTEMmv{b}  :   \textcolor{coeffectColor}{\square_{ \SYSTEMnt{r} } }   (   \SYSTEMnt{A}  \otimes  \SYSTEMnt{B}   )   }
        \inferrule*[Right=\SYSTEMRenameRuleLinprode{}]
          {
            \inferrule*[Right=\SYSTEMRenameRuleLinder{}*]
              {\inferrule*[Right=\SYSTEMRenameRuleLinvar{}]{ }{   \SYSTEMmv{p}  :   \SYSTEMnt{A}  \otimes  \SYSTEMnt{B}     \vdash_{\textsc{l} }  \SYSTEMmv{p}  :   \SYSTEMnt{A}  \otimes  \SYSTEMnt{B}  } \hspace{5em} r = 1 }
              {   \SYSTEMmv{p}  : \textcolor{coeffectColor}{[}   \SYSTEMnt{A}  \otimes  \SYSTEMnt{B}  {\textcolor{coeffectColor}{]_{ \SYSTEMnt{r} } } }    \vdash_{\textsc{l} }  \SYSTEMmv{p}  :   \SYSTEMnt{A}  \otimes  \SYSTEMnt{B}  }
            \hspace{4em}
            \inferrule*[Right=\SYSTEMRenameRuleLinprodi{}]
              {
                \inferrule*[Right=?]
                  { }
                  {    \SYSTEMmv{x}  :  \SYSTEMnt{A}    \vdash_{\textsc{l} }   \textcolor{coeffectColor}{[}  \SYSTEMmv{x}  \textcolor{coeffectColor}{]}   :   \textcolor{coeffectColor}{\square_{ \SYSTEMnt{r} } }  \SYSTEMnt{A}   }
                \hspace{4em}
                \inferrule*[Right=?]
                  { }
                  {    \SYSTEMmv{y}  :  \SYSTEMnt{B}    \vdash_{\textsc{l} }   \textcolor{coeffectColor}{[}  \SYSTEMmv{y}  \textcolor{coeffectColor}{]}   :   \textcolor{coeffectColor}{\square_{ \SYSTEMnt{r} } }  \SYSTEMnt{B}   }
              }
              {    \SYSTEMmv{x}  :   \textcolor{coeffectColor}{\square_{ \SYSTEMnt{r} } }  \SYSTEMnt{A}     \SYSTEMsym{+}    \SYSTEMmv{y}  :   \textcolor{coeffectColor}{\square_{ \SYSTEMnt{r} } }  \SYSTEMnt{B}     \vdash_{\textsc{l} }   \langle   \textcolor{coeffectColor}{[}  \SYSTEMmv{x}  \textcolor{coeffectColor}{]}  ,   \textcolor{coeffectColor}{[}  \SYSTEMmv{y}  \textcolor{coeffectColor}{]}   \rangle   :    \textcolor{coeffectColor}{\square_{ \SYSTEMnt{r} } }  \SYSTEMnt{A}   \otimes   \textcolor{coeffectColor}{\square_{ \SYSTEMnt{r} } }  \SYSTEMnt{B}   }
          }
          {   \SYSTEMmv{p}  : \textcolor{coeffectColor}{[}   \SYSTEMnt{A}  \otimes  \SYSTEMnt{B}  {\textcolor{coeffectColor}{]_{ \SYSTEMnt{r} } } }    \vdash_{\textsc{l} }   \mathsf{let} \, \langle  \SYSTEMmv{x} ,  \SYSTEMmv{y}  \rangle =  \SYSTEMmv{p}  \, \mathsf{in} \,   \langle   \textcolor{coeffectColor}{[}  \SYSTEMmv{x}  \textcolor{coeffectColor}{]}  ,   \textcolor{coeffectColor}{[}  \SYSTEMmv{y}  \textcolor{coeffectColor}{]}   \rangle    :    \textcolor{coeffectColor}{\square_{ \SYSTEMnt{r} } }  \SYSTEMnt{A}   \otimes   \textcolor{coeffectColor}{\square_{ \SYSTEMnt{r} } }  \SYSTEMnt{B}    }
      }
      {
        %\inferrule*[Right=\SYSTEMRenameRuleLinabs{}]
          {   \SYSTEMmv{b}  :   \textcolor{coeffectColor}{\square_{ \SYSTEMnt{r} } }   (   \SYSTEMnt{A}  \otimes  \SYSTEMnt{B}   )      \vdash_{\textsc{l} }   \mathsf{let} \, \textcolor{coeffectColor}{[}  \SYSTEMmv{p}  \textcolor{coeffectColor}{]} =  \SYSTEMmv{b}  \, \mathsf{in} \,   \mathsf{let} \, \langle  \SYSTEMmv{x} ,  \SYSTEMmv{y}  \rangle =  \SYSTEMmv{p}  \, \mathsf{in} \,   \langle   \textcolor{coeffectColor}{[}  \SYSTEMmv{x}  \textcolor{coeffectColor}{]}  ,   \textcolor{coeffectColor}{[}  \SYSTEMmv{y}  \textcolor{coeffectColor}{]}   \rangle     :    \textcolor{coeffectColor}{\square_{ \SYSTEMnt{r} } }  \SYSTEMnt{A}   \otimes   \textcolor{coeffectColor}{\square_{ \SYSTEMnt{r} } }  \SYSTEMnt{B}   }
        %  {  \emptyset   \vdash_{\textsc{l} }   \lambda  \SYSTEMmv{b}  .  \SYSTEMsym{(}   \mathsf{let} \, \textcolor{coeffectColor}{[}  \SYSTEMmv{p}  \textcolor{coeffectColor}{]} =  \SYSTEMmv{b}  \, \mathsf{in} \,   \mathsf{let} \, \langle  \SYSTEMmv{x} ,  \SYSTEMmv{y}  \rangle =  \SYSTEMmv{p}  \, \mathsf{in} \,   \langle   \textcolor{coeffectColor}{[}  \SYSTEMmv{x}  \textcolor{coeffectColor}{]}  ,   \textcolor{coeffectColor}{[}  \SYSTEMmv{y}  \textcolor{coeffectColor}{]}   \rangle     \SYSTEMsym{)}   :     \textcolor{coeffectColor}{\square_{ \SYSTEMnt{r} } }   (   \SYSTEMnt{A}  \otimes  \SYSTEMnt{B}   )    \multimap   \textcolor{coeffectColor}{\square_{ \SYSTEMnt{r} } }  \SYSTEMnt{A}    \otimes   \textcolor{coeffectColor}{\square_{ \SYSTEMnt{r} } }  \SYSTEMnt{B}   }
      }
\end{gather*}
%
At the points marked with question marks (?), we cannot apply promotion (\SYSTEMRenameRuleLinpr{}) to
$x$ and $y$, because the \textit{graded} predicate in the
antecedent does not hold; they are linear variables. Further, to apply \SYSTEMRenameRuleLinder{} at
the point marked with an asterisk (*), we would have to restrict to $r = 1$.
A similar attempted derivation for $  \textcolor{GLcolor}{\llbracket}  \SYSTEMmv{g}  \textcolor{GLcolor}{\rrbracket}  _\circ $, commuting the graded modality
and sum types, encounters the same issues.

One may wonder if any sound translation is possible at all. \citet{hughes2021linear} have explored
precisely this point and show that linear-base and graded-base systems exhibit a fundamental
difference in expressivity. Linear Core is akin to Linear Logic, in that it does
not admit derivation of the distributive law $!(A \otimes B) \multimap \, !A \otimes !B$ as the
semantics of multiplicative tensor $\otimes$ is that each component must be used equally; such a
distributive law would allow $A$ and $B$ to be decoupled and used differently. On the other hand,
Graded Base and the graded-base systems in the literature cited above admit such a distributive law
as seen with $f$. Thus Graded Base cannot model Linear Logic's $!$-modality exactly through \emph{any}
choice of semiring for the graded modality.

\citet{hughes2021linear} propose extensions to graded systems to bring back the restrictions of
Linear Logic. Here we explore a different route, namely we extend Linear Core by further constructs
which admit the augmented expressivity of Graded Modal Core, giving the system \textbf{Linear~Push~Core}.

%We could `drag back' the expressivity of Graded Base by requiring that
%$r = 1$ in the product and sum eliminations, however even that is
%still more general.

\subsection{Linear Push Core}
\label{sec:linear-push-core}
%
We extend Linear Core to \textbf{Linear~Push~Core} with new syntax, heaving it to a level of
expressivity on par with Graded Modal Core, leaving behind the restrictions of Linear Logic
discussed in \ref{sec:no-trans-graded-core-to-linear-core}.

Since the essential missing piece is the distributive laws between the graded modality and
products and sums, we complete the syntax of Linear Push Core with the following two constructs:
%
\begin{align*}
  \SYSTEMnt{t} & ::= ... \mid  \textsf{push}_\otimes  \SYSTEMnt{t}  \mid  \textsf{push}_\oplus  \SYSTEMnt{t}  \mid  \textsf{push}_{\mathrm{unit} }  \SYSTEMnt{t} 
  \tag{terms}
\end{align*}
%
The syntax is non-standard but isomorphic to other approaches; Granule achieves the same
expressivity implicitly via pattern matching~\cite{DBLP:journals/pacmpl/OrchardLE19}, but
we prefer to introduce explicit syntax here. These two constructs provide exactly
the distributive behaviour captured in product and sum elimination in Graded Modal Core.
Figure~\ref{fig:linear-core-push} gives the typing and operational model.

\begin{figure}[t]
\raggedright\noindent\framebox{$ \Gamma  \vdash_{\textsc{l} }  \SYSTEMnt{t}  :  \SYSTEMnt{A} $} \quad Typing rules
\begin{gather*}
\begin{align*}
\begin{array}{c}
\SYSTEMdruleLinpushprod{}
\,
\SYSTEMdruleLinpushsum{}
\,
\SYSTEMdruleLinpushunit{}
\end{array}
\end{align*}
\end{gather*}
\bigskip

\raggedright\noindent\framebox{$ \SYSTEMnt{t}  \rightsquigarrow_{\textsc{l} }  \SYSTEMnt{t'} $} \quad Operational semantics for push constructors
\begin{gather*}
\begin{align*}
\begin{array}{c}
    \SYSTEMdruleSemLinpushProd{}
   \;\;
   \SYSTEMdruleSemLinpushUnit{}
   \;\;
   \SYSTEMdruleSemLinpushSumInjGen{}
  \\[1em]
  \SYSTEMdruleSemLinpushProdCong{}
  \;
  \SYSTEMdruleSemLinpushSumCong{}
  \;
  \SYSTEMdruleSemLinpushUnitCong{}
  \\[1em]
  \SYSTEMdruleSemLinpushProdBoxCong{}
  \;
  \SYSTEMdruleSemLinpushSumBoxCong{}
  \;
  \SYSTEMdruleSemLinpushUnitBoxCong{}
  \end{array}
\end{align*}
\end{gather*}
\caption{Linear Push Core}
\label{fig:linear-core-push}
\end{figure}

\subsection{Translation from Graded Modal Core to Linear Push Core}
\label{sec:trans-graded-core-to-linear-push-core}
%
We define the translation
from Graded Modal Core into Linear Push Core as follows
(fixing the unsound translation to Linear Core from Section~\ref{sec:no-trans-graded-core-to-linear-core}
in the product and sum elimination):
%
%
\renewcommand{\triangleq}{& \oldTriangleq &}
\renewcommand{\equiv}{& \oldEquiv &}
\renewcommand{\SYSTEMdrule}[4][]{#3 \qquad #2}
%
\begin{gather*}
{\setlength{\arraycolsep}{0.1em}
\begin{array}{lcl@{\hspace{-4em}}r}
\SYSTEMdruleTranslGtoLTmprodi{} \\
 \textcolor{GLcolor}{\llbracket}   \mathsf{let} \, \langle  \SYSTEMmv{x} ,  \SYSTEMmv{y}  \rangle =  \SYSTEMnt{t_{{\mathrm{1}}}}  \, \mathsf{in} \,  \SYSTEMnt{t_{{\mathrm{2}}}}   \textcolor{GLcolor}{\rrbracket}  \triangleq  \mathsf{let} \, \langle  \SYSTEMmv{x'} ,  \SYSTEMmv{y'}  \rangle =   \textsf{push}_\otimes   \textcolor{coeffectColor}{[}   \textcolor{GLcolor}{\llbracket}  \SYSTEMnt{t_{{\mathrm{1}}}}  \textcolor{GLcolor}{\rrbracket}   \textcolor{coeffectColor}{]}    \, \mathsf{in} \,   \mathsf{let} \, \textcolor{coeffectColor}{[}  \SYSTEMmv{x}  \textcolor{coeffectColor}{]} =  \SYSTEMmv{x'}  \, \mathsf{in} \,   \mathsf{let} \, \textcolor{coeffectColor}{[}  \SYSTEMmv{y}  \textcolor{coeffectColor}{]} =  \SYSTEMmv{y'}  \, \mathsf{in} \,   \textcolor{GLcolor}{\llbracket}  \SYSTEMnt{t_{{\mathrm{2}}}}  \textcolor{GLcolor}{\rrbracket}     \\
\SYSTEMdruleTranslGtoLTmsumiOne{} \\
\SYSTEMdruleTranslGtoLTmsumiTwo{} \\
  \textcolor{GLcolor}{\llbracket}   \mathsf{case} \,  \SYSTEMnt{t}  \, \mathsf{of} \, \{ \mathsf{inj1} \,  \SYSTEMmv{x}  \rightarrow  \SYSTEMnt{t_{{\mathrm{1}}}}  ; \, \mathsf{inj2} \,  \SYSTEMmv{y}  \rightarrow  \SYSTEMnt{t_{{\mathrm{2}}}}  \}   \textcolor{GLcolor}{\rrbracket}  \\
\multicolumn{3}{l}{\hspace{0.2em}\;\quad \oldTriangleq \;  \mathsf{case} \,   \textsf{push}_\oplus   \textcolor{coeffectColor}{[}   \textcolor{GLcolor}{\llbracket}  \SYSTEMnt{t}  \textcolor{GLcolor}{\rrbracket}   \textcolor{coeffectColor}{]}    \, \mathsf{of} \, \{ \mathsf{inj1} \,  \SYSTEMmv{x'}  \rightarrow   \mathsf{let} \, \textcolor{coeffectColor}{[}  \SYSTEMmv{x}  \textcolor{coeffectColor}{]} =  \SYSTEMmv{x'}  \, \mathsf{in} \,   \textcolor{GLcolor}{\llbracket}  \SYSTEMnt{t_{{\mathrm{1}}}}  \textcolor{GLcolor}{\rrbracket}    ; \, \mathsf{inj2} \,  \SYSTEMmv{y'}  \rightarrow   \mathsf{let} \, \textcolor{coeffectColor}{[}  \SYSTEMmv{y}  \textcolor{coeffectColor}{]} =  \SYSTEMmv{y'}  \, \mathsf{in} \,   \textcolor{GLcolor}{\llbracket}  \SYSTEMnt{t_{{\mathrm{2}}}}  \textcolor{GLcolor}{\rrbracket}    \}  }\\
\SYSTEMdruleTranslGtoLTmuniti{} \\
\SYSTEMdruleTranslGtoLTmunite{} & \text{(Term translation)} \\[0.5em]
\SYSTEMdruleTranslGtoLTyprodTy{} \\
\SYSTEMdruleTranslGtoLTysumTy{} \\
\SYSTEMdruleTranslGtoLTyunitTy{} & \text{(Type translation)}
\end{array}}
\end{gather*}
%
\let\SYSTEMdrule\oldSYSTEMdrule
\let\equiv\oldEquiv
\let\triangleq\oldTriangleq
%
Translation of a product elimination on $\SYSTEMnt{t_{{\mathrm{1}}}}$ first promotes
the product to a graded modality which is then `pushed into' each component of the product via $\mathsf{push}_\otimes$. The resulting graded modal components $x'$ and $y'$ are then eliminated in the scope of $\SYSTEMnt{t_{{\mathrm{2}}}}$. A similar idea is applied for sum elimination. Thus this translation
gets at the heart of what is happening inside Graded Modal Core (and why it differs from
Linear Core): its eliminators distribute the coeffect analysis inside products and sums,
which we now capture explicitly here in Linear Push Core.

We now re-establish the core results (proofs for which are
in Appendix~\refappendix{app:proofs-grad-core-to-lin-core}):
%
\begin{restatable}{thrm}{gradToLinTranslationExt}[Soundness of the Graded Modal Core to Linear Push Core translation]
 \begin{itemize}[itemsep=0.2em]
  \item Type preservation: $ \Delta  \vdash_{\textsc{g} }  \SYSTEMnt{t}  :  \SYSTEMnt{A}  \implies   \textcolor{GLcolor}{\llbracket}  \Delta  \textcolor{GLcolor}{\rrbracket}   \vdash_{\textsc{l} }   \textcolor{GLcolor}{\llbracket}  \SYSTEMnt{t}  \textcolor{GLcolor}{\rrbracket}   :   \textcolor{GLcolor}{\llbracket}  \SYSTEMnt{A}  \textcolor{GLcolor}{\rrbracket}  $;
  \item Operational correspondence: $ \SYSTEMnt{t}  \rightsquigarrow_{\textsc{g} }  \SYSTEMnt{t'}  \implies   \textcolor{GLcolor}{\llbracket}  \SYSTEMnt{t}  \textcolor{GLcolor}{\rrbracket}   \rightsquigarrow_{\textsc{l} }^\ast   \textcolor{GLcolor}{\llbracket}  \SYSTEMnt{t'}  \textcolor{GLcolor}{\rrbracket}  $.
  \item Equation preservation: $(\SYSTEMnt{t}  \equiv^{-\eta\otimes\oplus\mathsf{unit} }_{\textsc{g} }  \SYSTEMnt{t'}) \implies   \textcolor{GLcolor}{\llbracket}  \SYSTEMnt{t}  \textcolor{GLcolor}{\rrbracket}   \,\equiv_{\textsc{l} }\,   \textcolor{GLcolor}{\llbracket}  \SYSTEMnt{t'}  \textcolor{GLcolor}{\rrbracket}  $.
\end{itemize}
Our translation does not, in general, preserve $\eta$-equalities for
products, sums, and unit (it can only preserve $\eta$ on these types
for closed terms); i.e., our translation is not extensional due to
additional `push'-structure for handling nesting of graded modalities
and products, sums, and units.
\end{restatable}
%
Given that we have changed Linear Core to Linear Push Core by adding two new constructs,
we now re-establish that we can translate Linear Push Core back into Graded Modal Core.

\subsection{Translation from Linear Push Core to Graded Modal Core}
\label{sec:trans-linear-push-core-to-graded-core}
%
We extend the Linear Core to Graded Modal Core translation with the following
additional cases to handle the extra constructs of Linear Push Core:
%
\renewcommand{\equiv}{ \oldEquiv}
\renewcommand{\SYSTEMdrule}[4][]{{\displaystyle{#3}}}
%
\begin{align}
\tag{Term translation} \\[-1.5em]
\notag & \begin{array}{l}
\SYSTEMdruleTranslLtoGTmpushProd{} \\
\SYSTEMdruleTranslLtoGTmpushSum{} \\
\SYSTEMdruleTranslLtoGTmpushUnit{}
\end{array}
\end{align}
%
\let\SYSTEMdrule\oldSYSTEMdrule
\let\equiv\oldEquiv
%
Note that (for $ \textsf{push}_\otimes     $ and $ \textsf{push}_\oplus     $) these are essentially the programs $f$ and $g$ shown in Section~\ref{sec:no-trans-graded-core-to-linear-core}
which are the derivation of $\textsf{push}$ combinators in Graded Modal Core.

We thus (re)establish the key results (proofs in Appendix~\refappendix{app:proofs-lin-core-to-grad-core}):


\begin{restatable}{thrm}{linToGradTranslationExt}[Soundness of the Linear Push Core to Graded Modal Core translation]
 \begin{itemize}[itemsep=0.2em]
  \item Type preservation: $ \Gamma  \vdash_{\textsc{l} }  \SYSTEMnt{t}  :  \SYSTEMnt{A}  \implies   \textcolor{LGcolor}{\llparenthesis}  \Gamma  \textcolor{LGcolor}{\rrparenthesis}   \vdash_{\textsc{g} }   \textcolor{LGcolor}{\llparenthesis} \smidge  \SYSTEMnt{t}  \smidge \textcolor{LGcolor}{\rrparenthesis}   :   \textcolor{LGcolor}{\llparenthesis} \smidge  \SYSTEMnt{A}  \smidge \textcolor{LGcolor}{\rrparenthesis}  $;
  \item Operational correspondence: $ \SYSTEMnt{t}  \rightsquigarrow_{\textsc{l} }  \SYSTEMnt{t'}  \implies   \textcolor{LGcolor}{\llparenthesis} \smidge  \SYSTEMnt{t}  \smidge \textcolor{LGcolor}{\rrparenthesis}   \rightsquigarrow_{\textsc{g} }^\ast   \textcolor{LGcolor}{\llparenthesis} \smidge  \SYSTEMnt{t'}  \smidge \textcolor{LGcolor}{\rrparenthesis}  $;
  \item Equation preservation: $ \SYSTEMnt{t}  \,\equiv_{\textsc{l} }\,  \SYSTEMnt{t'}  \implies   \textcolor{LGcolor}{\llparenthesis} \smidge  \SYSTEMnt{t}  \smidge \textcolor{LGcolor}{\rrparenthesis}   \equiv_{\textsc{g} }   \textcolor{LGcolor}{\llparenthesis} \smidge  \SYSTEMnt{t'}  \smidge \textcolor{LGcolor}{\rrparenthesis}  $.
\end{itemize}
\end{restatable}
%
Note that whilst the operational correspondence from Linear Base to Graded Modal Base mapped a single
reduction $ \SYSTEMnt{t}  \rightsquigarrow_{\textsc{l} }  \SYSTEMnt{t'} $ to a single graded-base reduction $  \textcolor{LGcolor}{\llparenthesis} \smidge  \SYSTEMnt{t}  \smidge \textcolor{LGcolor}{\rrparenthesis}   \rightsquigarrow_{\textsc{g} }   \textcolor{LGcolor}{\llparenthesis} \smidge  \SYSTEMnt{t'}  \smidge \textcolor{LGcolor}{\rrparenthesis}  $ this
is now not possible for Linear Push Core. Instead, a single reduction in Linear Push Core for
the $\mathsf{push}$ combinators corresponds to several reductions in Graded Modal Core.


\iffalse
% SOME LEFTOVERS HERE
\begin{align*}
 \textcolor{GLcolor}{\llbracket}  \SYSTEMmv{f}  \textcolor{GLcolor}{\rrbracket}  &:   \textcolor{coeffectColor}{\square_{ \SYSTEMnt{r} } }   (   \SYSTEMnt{A}  \otimes  \SYSTEMnt{B}   )    \multimap   (    \textcolor{coeffectColor}{\square_{ \SYSTEMnt{r} } }  \SYSTEMnt{A}   \otimes   \textcolor{coeffectColor}{\square_{ \SYSTEMnt{r} } }  \SYSTEMnt{B}    )   \\
 \textcolor{GLcolor}{\llbracket}  \SYSTEMmv{f}  \textcolor{GLcolor}{\rrbracket}  &=  \lambda  \SYSTEMmv{b}  .  \SYSTEMsym{(}   \mathsf{let} \, \textcolor{coeffectColor}{[}  \SYSTEMmv{p}  \textcolor{coeffectColor}{]} =  \SYSTEMmv{b}  \, \mathsf{in} \,   \mathsf{let} \, \langle  \SYSTEMmv{x'} ,  \SYSTEMmv{y'}  \rangle =   \textsf{push}_\otimes   \textcolor{coeffectColor}{[}  \SYSTEMmv{p}  \textcolor{coeffectColor}{]}    \, \mathsf{in} \,   \mathsf{let} \, \textcolor{coeffectColor}{[}  \SYSTEMmv{x}  \textcolor{coeffectColor}{]} =  \SYSTEMmv{x'}  \, \mathsf{in} \,   \mathsf{let} \, \textcolor{coeffectColor}{[}  \SYSTEMmv{y}  \textcolor{coeffectColor}{]} =  \SYSTEMmv{y'}  \, \mathsf{in} \,   \langle   \textcolor{coeffectColor}{[}  \SYSTEMmv{x}  \textcolor{coeffectColor}{]}  ,   \textcolor{coeffectColor}{[}  \SYSTEMmv{y}  \textcolor{coeffectColor}{]}   \rangle       \SYSTEMsym{)} \\
\\
 \textcolor{GLcolor}{\llbracket}  \SYSTEMmv{g}  \textcolor{GLcolor}{\rrbracket}  &:   \textcolor{coeffectColor}{\square_{ \SYSTEMnt{r} } }   (   \SYSTEMnt{A}  \oplus  \SYSTEMnt{B}   )    \multimap   (    \textcolor{coeffectColor}{\square_{ \SYSTEMnt{r} } }  \SYSTEMnt{A}   \oplus   \textcolor{coeffectColor}{\square_{ \SYSTEMnt{r} } }  \SYSTEMnt{B}    )   \\
 \textcolor{GLcolor}{\llbracket}  \SYSTEMmv{g}  \textcolor{GLcolor}{\rrbracket}  &=  \lambda  \SYSTEMmv{b}  .  \SYSTEMsym{(}   \mathsf{let} \, \textcolor{coeffectColor}{[}  \SYSTEMmv{e}  \textcolor{coeffectColor}{]} =  \SYSTEMmv{b}  \, \mathsf{in} \,   \mathsf{case} \,   \textsf{push}_\oplus   \textcolor{coeffectColor}{[}  \SYSTEMmv{e}  \textcolor{coeffectColor}{]}    \, \mathsf{of} \, \{ \mathsf{inj1} \,  \SYSTEMmv{x'}  \rightarrow   \mathsf{let} \, \textcolor{coeffectColor}{[}  \SYSTEMmv{x}  \textcolor{coeffectColor}{]} =  \SYSTEMmv{x'}  \, \mathsf{in} \,   \mathsf{inj}_1 \,   \textcolor{coeffectColor}{[}  \SYSTEMmv{x}  \textcolor{coeffectColor}{]}     ; \, \mathsf{inj2} \,  \SYSTEMmv{y}  \rightarrow   \mathsf{inj}_2 \,   \textcolor{coeffectColor}{[}  \SYSTEMmv{y}  \textcolor{coeffectColor}{]}    \}    \SYSTEMsym{)} 
\end{align*}
\fi

\begin{remark}
Graded Modal Core extends Graded Modal Base rather than Graded Base.
If Graded Modal Core did not have a graded modality it could not give a precise type to projection on products:
%
\begin{align*}
  \emptyset   \vdash_{\textsc{g} }  \SYSTEMsym{(}   \lambda  \SYSTEMmv{z}  .   \mathsf{let} \, \langle  \SYSTEMmv{x} ,  \SYSTEMmv{y}  \rangle =  \SYSTEMmv{z}  \, \mathsf{in} \,  \SYSTEMmv{x}    \SYSTEMsym{)}  :    (   \SYSTEMnt{A}  \times  \SYSTEMnt{B}   )   \xrightarrow{\textcolor{coeffectColor}{ \SYSTEMnt{r} } }  \SYSTEMnt{A}  
\end{align*}
%
where $ \SYSTEMsym{1}  \, \textcolor{coeffectColor}{\sqsubseteq} \,  \SYSTEMnt{r} $ due to the use of $x$ and $ \SYSTEMsym{0}  \, \textcolor{coeffectColor}{\sqsubseteq} \,  \SYSTEMnt{r} $ due to the
absence of $y$.  For the exact-usage natural numbers semiring this
would not be typeable as there is no such $r$ satisfying both
constraints. For a semiring of intervals denoting lower-bound and
upper-bound usage~\cite{DBLP:journals/pacmpl/OrchardLE19}, the
constraints would be satisfied by $r = [0,1]$.

We instead opted to include the graded modality in Graded Modal Core as in the work of \citet{DBLP:journals/pacmpl/AbelB20}, \citet{grtt}, and others~\cite{DBLP:conf/esop/HughesO24,mycroftfest2024}.
Thus we can specify for this example a more fine-grained type via the graded modality,
where the $\SYSTEMnt{B}$ component of the product is unused:
%
\[
    \emptyset   \vdash_{\textsc{g} }   \lambda  \SYSTEMmv{z}  .   \mathsf{let} \, \langle  \SYSTEMmv{x} ,  \SYSTEMmv{y}  \rangle =  \SYSTEMmv{z}  \, \mathsf{in} \,   \mathsf{let} \, \textcolor{coeffectColor}{[}  \SYSTEMmv{y'}  \textcolor{coeffectColor}{]} =  \SYSTEMmv{y}  \, \mathsf{in} \,  \SYSTEMmv{x}     :    (   \SYSTEMnt{A}  \times   \textcolor{coeffectColor}{\square_{ \SYSTEMsym{0} } }  \SYSTEMnt{B}    )   \xrightarrow{\textcolor{coeffectColor}{ \SYSTEMsym{1} } }  \SYSTEMnt{A}  
\]
%
Without the inclusion of graded modalities, Graded Modal Core is less expressive than Linear Core.

Note that some dependently-typed graded-base systems such as \textsc{GraD}~\cite{grad}
and QTT~\cite{quantitative-type-theory} omit the graded modality but include a graded
dependent $\Sigma$ type which can be used to encode the graded modality.
\end{remark}
